\begin{frame}{좋은 리뷰를 쓰는 과정: 체크리스트 $\to$ 반박 가능한 질문}
  \begin{columns}[T]
    \begin{column}{0.5\textwidth}
      \centering
      % [fig missing] \includegraphics[width=\linewidth]{ch08_2_rl_review_flow.png}
      \vskip0.3em
      {\footnotesize 발표+보고서로 주장 파악 $\to$ 4항목 적용 $\to$
      각 항목에 질문 초안(답이 \textbf{새로운 계산·실험}이어야 함)
      $\to$ ``개념(장·절)을 근거로 들고, 답이 새 실험인가?'' 분기: 예 $\to$
      \textbf{채택}, 아니오 $\to$ 감상평으로 떨어지고 다시 쓰기 $\to$
      3단계(1/2/3점)로 자기 채점.}
    \end{column}
    \begin{column}{0.48\textwidth}
      \begin{itemize}
        \item 흐름: 발표를 듣고 $\to$ 4항목에 적용 $\to$ 각 항목에서
          \textbf{구체적 관찰}을 적고 $\to$ 이 학기 개념(장·절)에
          뿌리를 박은 질문인지 검사 $\to$ 아니면 \textbf{다시 쓰기}.
        \item RL 버전의 특징: ML1 템플릿에서 \textbf{5번을 두 칸으로
          늘림} --- ``근거가 되는 \textbf{개념}'' + ``반박하려면 어떤
          \textbf{새 실험}이 필요한지''.
        \item 답이 \textbf{새 실험}이어야만 질문이 3점.
      \end{itemize}
    \end{column}
  \end{columns}
\end{frame}
